\documentclass[10pt,a4paper]{article}
\usepackage[margin=25mm]{geometry}
\usepackage{amsmath,amssymb,amsthm,hyperref}
\setlength{\emergencystretch}{2em}
\newtheorem{lemma}{Lemma}
\newtheorem{theorem}[lemma]{Theorem}
\newcommand{\gen}[1]{\langle#1\rangle}
\title{Natural paths on the residual fixed points}
\author{Human Compression project: L044 continuation}
\date{}
\begin{document}\maketitle
Let $G=\gen{\rho_0,\ldots,\rho_{r-1}}\leq A_n$ be a transitive
connected string C-group, with every maximal omission intransitive.
Assume $G_0$ has orbits $\{a\}$ and its complement, and put
\[
 \rho_0=(a\ b)\delta,\quad X=\Omega\setminus\{a,b\},\quad
 K=G_{>1},\quad R=G_{>2},\quad \sigma=\rho_1.
\]
The odd involution $\delta$ has $t$ pairs, with set $\mathcal T$
and union $P$, and fixed set $F\subset X$ of size $f$.
Thus $K$ fixes $a,b$, commutes with $\delta$, and preserves $F,P$.
The involution $\sigma$ fixes $a$ and centralizes $R$.
Throughout assume $n=2r\geq14$, so $t$ is odd, $f$ is even, and
the number of distinguished $K$ labels is $d=t+f/2-1$.
We use the complete accepted actual block-orbit deficit proof:
every distinguished image on the $\delta$-blocks is orbit-essential.
Whiston's bound on independent generators of subgroups of $S_m$
remains a retained input, with its complete paid predecessor proof interfaces.

\begin{lemma}[Transport and the pair budget]\label{transport}
Let $U\subseteq F$ be an $R$-orbit whose induced action has trivial
centralizer in $\operatorname{Sym}(U)$. If $O=\sigma(U)\subseteq P$,
then $O$ and $\delta(O)$ are disjoint $R$-orbits and use exactly
$|U|$ pairs of $\mathcal T$. Any subfamily of distinguished $R$
generators fixing $U$ pointwise also fixes these pairs pointwise.
If its pair images are independent, it has at most
$\max\{0,t-|U|-1\}$ members.
\end{lemma}
\begin{proof}
Both permutations commute with $R$, so transport orbits equivariantly.
If $\delta(O)=O$, its restriction centralizes the transported action
and is trivial, contradicting $O\subset P$. Hence the orbits are
disjoint, paired by $\delta$. Equivariance transfers pointwise fixation
from $U$ to both copies. The independent pair images act on the remaining
$t-|U|$ pairs; Whiston gives the stated bound, with zero for a set
of size zero or one. In particular a natural symmetric action of
degree at least three satisfies the centralizer hypothesis.
\end{proof}

\begin{lemma}[Routing a natural fixed-set path]\label{route}
Suppose $f\geq4$, and the nonidentity distinguished restrictions to
$F$ are exactly the $f-1$ successive transpositions of a natural
$S_f$ path. All other distinguished labels, denoted $M$, fix $F$
pointwise. Then the fixed-set path occupies indices $2,\ldots,f$.
Number its vertices $z,u_1,\ldots,u_{f-1}$ in this order.
The permutation $\sigma$ fixes $U=\{u_1,\ldots,u_{f-1}\}$ pointwise
and sends $z$ to a point $x\in P$. The whole $R$ fixes the pair
$C=\{x,x^\delta\}$ pointwise.

There are $q=t-f/2$ labels in $M$, with independent and orbit-essential
pair images, and they all fix $C$. Orbit-essential here refers to
the action of $M$ itself, not only to the larger $K$ action.
\end{lemma}
\begin{proof}
Noncommuting consecutive path transpositions force their indices to
be consecutive in the full string. The count gives $|M|=q$.
Every $M$ pair image is independent of the others: a dependence
would also be a dependence of full block images, since $M$ fixes $F$.
More strongly, if deleting an $M$ label left the $M$ pair-orbits
unchanged, adjoining all fixed-set path labels would leave the full
$K$ block-orbits unchanged. This contradicts the accepted actual-orbit
deficit. Thus these pair images are orbit-essential.

If $\rho_2$ belonged to $M$, then $R$ would still induce natural
$S_f$ on $F$. Transitivity of $G_0$ forbids $\sigma(F)=F$, so
$\sigma(F)$ would be an $R$-orbit in $P$. Lemma~\ref{transport}
would force $t\geq f$ and give a bound $t-f-1$ when $t>f$,
or zero when $t=f$, for the $q-1=t-f/2-1$ labels of
$M\setminus\{\rho_2\}$. Both are contradictions for even $f\geq4$.
Therefore the path begins at index2.

Now $R$ has its natural $S_{f-1}$ orbit $U$ and fixes $z$.
If $\sigma(U)\ne U$, its image must be an orbit in $P$.
The transport lemma gives $t\geq f-1$ and bounds all $q$ labels
of $M$ by $t-f$ if $t\geq f$, and by zero if $t=f-1$.
Both contradict $q=t-f/2$. Thus $\sigma$ preserves $U$;
its centralizer is trivial, so $\sigma$ fixes it pointwise.

The point $\sigma(z)$ cannot lie in $U$ or equal $a$.
If it were $z$, then $G_0$ would preserve $F$; if it were $b$,
then, as $\sigma$ is an involution, $G_0$ would preserve
$F\cup\{b\}$. Both are proper subsets of the complementary
$G_0$ orbit, because $P$ is nonempty. Hence $x=\sigma(z)\in P$.
Centralization and the fact that $R$ fixes $z$ show that $R$ fixes
$x$; commutation with $\delta$ also fixes $x^\delta$.
All labels in $M$ lie in $R$, proving the final assertion.
\end{proof}

\begin{theorem}[The six-point natural-path case]\label{six}
The hypotheses of Lemma~\ref{route} are impossible when $f=6$.
\end{theorem}
\begin{proof}
Here $t\geq3$ is odd. The fixed-set path has labels $2,3,4,5,6$;
$M$ consists of labels at least7. Put $m=t-1$, and consider its
pair action on $\mathcal T\setminus\{C\}$. It has $q=t-3=m-2$
orbit-essential distinguished generators, where the empty case
$m=2,q=0$ is allowed. Each of $\rho_3,\rho_4,\rho_5$ commutes with
$M$ and $\delta$, fixes $C$ pointwise, and is odd on the remaining
paired points: it is a transposition on $F$ and even globally.

The $M$ pair action cannot be transitive. Otherwise
$\gen{M,\delta}$ is transitive on its $2m$ points; a commuting
involution is either trivial or fixed-point-free, and hence has
either zero or $m$ transpositions. Both numbers are even, a
contradiction. There are therefore at least two pair orbits.
The orbit-essential forest has $m$ vertices and $m-2$ edges,
so there are exactly two, say $A,B$, of sizes $\alpha,\beta$.
The forest spans each of these two orbits.

Saturation determines the pair images exactly. For any selected
forest edge $e$, all generators other than its label preserve the
two pieces of the cut obtained by deleting $e$. For a given label,
intersect the cuts of every other edge. All resulting classes are
singletons except for the two endpoints of its own edge. Its pair
image is therefore precisely their transposition. Consequently
the pair image of $M$ is the full $S_A\times S_B$.
This includes a singleton orbit and, when $m=2$, two singleton orbits.

Let $\xi$ be any of $\rho_3,\rho_4,\rho_5$. Its pair image centralizes
$S_A\times S_B$. Unless $\alpha=\beta=1$, it preserves each of $A,B$:
the nontrivial natural symmetric factors are independent, and their
fixed points determine a possible singleton factor. If both sizes
were even, $\xi$ would be even on the points over each orbit.
Indeed it centralizes the transitive action of
$\gen{M|_{P_A},\delta|_{P_A}}$ there, so its restriction has zero
or $\alpha$ transpositions, and similarly on $P_B$.
Since $\alpha+\beta=m$ is even and $\xi$ is odd on these points,
both sizes must be odd.

For an odd size at least three the natural symmetric centralizer
is trivial; a singleton is also fixed. Hence the pair image of
$\xi$ is trivial, and its flip vector is constant on each of
$A,B$. Oddness makes it flip exactly one of the two odd-sized
orbits. In the remaining case $\alpha=\beta=1$, an involution
interchanging the two pair-blocks has two point transpositions,
and is even. Thus oddness again forces the pair image to be trivial
and exactly one of the two pairs to be flipped.

In all cases the pair part of $\rho_4$ is the full flip on a
nonempty proper subset $W=P_A$ or $W=P_B$. Every distinguished
$K$ generator other than $\rho_3,\rho_5$ commutes with $\rho_4$;
because it preserves $F$ and $P$, it preserves this subset $W$.
The two exceptional generators are themselves flips on $P_A$ or
$P_B$, and preserve $W$ as well. Hence $K$ preserves $W$.
Finally $\sigma$ commutes with $\rho_4$ and fixes $U$ pointwise.
The two points moved by $\rho_4$ in $F$ lie in $U$ and are fixed
by $\sigma$, so $\sigma$ must preserve the rest of its support,
namely $W$. Thus $G_0=\gen{\sigma,K}$ preserves $W$, contradicting
its transitivity on $B$. This proves the theorem.
\end{proof}

\paragraph{Attribution and boundary.}
The source is Cameron--Fernandes--Leemans--Mixer,
\emph{Highest rank of a polytope for $A_n$}, arXiv:1605.09173v1,
especially its singleton-end strategy on physical pages38--44.
The natural-orbit transport argument was introduced in the separate
project proof of the transitive four-point case by the collaborator
\texttt{rewrite\_termination}. Lemmas~\ref{transport}--\ref{route}
record its general-degree form; the two-orbit saturation and
middle-label support obstruction in Theorem~\ref{six} are the present
continuation. The preceding actual block-deficit proof remains paid.

This does not classify arbitrary actions on six fixed points.
Its five literal path restrictions are hypotheses, not consequences
of transitivity or of a quotient intersection property. The routing
lemma alone leaves all other even fixed-set sizes open. There is no
smaller-degree alternating induction, new finite computation or formal
claim here, and no whole L044, source-credit or question-credit claim.
Distinct review of these statements remains pending.
\end{document}
